\documentclass[11pt]{article}

% Preamble
\usepackage{amsmath,amssymb,amsthm}
\usepackage{natbib}
\usepackage{hyperref}
\usepackage{graphicx}
\usepackage{booktabs}
\usepackage{algorithm}
\usepackage{algorithmic}
\usepackage[margin=1in]{geometry}

% Theorem environments
\newtheorem{theorem}{Theorem}[section]
\newtheorem{proposition}[theorem]{Proposition}
\newtheorem{lemma}[theorem]{Lemma}
\newtheorem{corollary}[theorem]{Corollary}
\newtheorem{definition}[theorem]{Definition}
\newtheorem{assumption}{Assumption}
\newtheorem{remark}{Remark}

% Custom commands
\newcommand{\R}{\mathbb{R}}
\newcommand{\E}{\mathbb{E}}
\newcommand{\Prob}{\mathbb{P}}
\newcommand{\Var}{\mathrm{Var}}
\newcommand{\Cov}{\mathrm{Cov}}
\newcommand{\ind}{\mathbf{1}}

\title{Extreme Conformal Prediction for Intraday Tail Events}
\author{Author}
\date{\today}

\begin{document}

\maketitle

\begin{abstract}
Conformal prediction provides distribution-free coverage guarantees, but at extreme quantiles (99.9\%+) classical methods yield uninformative infinite-width intervals. Financial regulators mandate risk measures at precisely these extreme levels. We bridge extreme value theory and conformal prediction under temporal dependence, using Generalized Pareto Distribution (GPD) extrapolation to extend conformal calibration beyond empirical quantiles while the Barber-Pananjady switch coefficient framework provides coverage guarantees under $\beta$-mixing. Our coverage bound decomposes as $O(\beta(\tau) + n_u^{-1/2})$ where $\beta(\tau)$ captures mixing-induced degradation and $n_u$ is the number of threshold exceedances used for GPD fitting. We introduce a censoring correction for circuit breaker effects that truncate observed extremes. An impossibility result establishes that our method targets heteroscedastic tails specifically: under homoscedastic tail behavior, GPD extrapolation cannot improve on empirical quantiles. Empirical validation on FX flash crashes (SNB 2015, Brexit 2016, JPY 2019) and cryptocurrency tail events (COVID 2020, FTX 2022) shows 30--50\% narrower intervals than classical conformal methods while maintaining nominal coverage.
\end{abstract}

% Sections
%!TEX root = ../main.tex

\section{Introduction}
\label{sec:introduction}

% Paragraph 1: Problem statement — extreme quantiles and regulatory requirements
Financial regulators mandate risk measures at extreme quantiles. Basel III requires Value-at-Risk at the 99th percentile for market risk capital, and Expected Shortfall at 97.5\% for internal models. Stress testing scenarios demand coverage guarantees at 99.9\% and beyond—the once-per-decade tail events that determine institutional solvency. Conformal prediction offers distribution-free, finite-sample coverage guarantees for prediction intervals: under exchangeability, a conformalized interval contains the true response with probability at least $1-\alpha$. Yet at extreme confidence levels, classical conformal methods fail. When the desired $\alpha$ is smaller than $1/(n+1)$, where $n$ is the calibration set size, the empirical quantile of conformity scores is undefined. The conformal interval becomes the entire real line—valid coverage, but operationally useless. A calibration set of 1,000 observations cannot distinguish 99.9\% from 100\% confidence via empirical quantiles. The problem compounds for high-frequency financial data: flash crashes, circuit breaker halts, and liquidity crises are precisely the tail events that matter for risk management, yet their rarity makes empirical calibration infeasible.

% Paragraph 2: Literature position and exact gap
Two research threads address parts of this problem. Extreme value theory (EVT) provides principled extrapolation beyond the data range: the Pickands-Balkema-de Haan theorem establishes that exceedances above a high threshold follow a Generalized Pareto Distribution (GPD), enabling estimation of arbitrary tail quantiles from moderate samples. \citet{pasche2025extreme} recently bridged EVT and conformal prediction for environmental extremes, using GPD extrapolation to construct conservative prediction intervals at confidence levels well beyond what empirical quantiles can reach. Their approach assumes approximate exchangeability—appropriate for flood forecasting with seasonal adjustment, but violated by the serial dependence inherent in financial returns. Separately, \citet{barber2025predictive} provide conformal coverage guarantees under temporal dependence via the \emph{switch coefficient} framework: for $\beta$-mixing processes, coverage degradation scales with a dependence measure $\bar{\Psi}_\tau$ that captures how far blocks of scores depart from independence. Their bounds apply to any conformity score but do not address the infinite-interval problem at extreme quantiles. The gap we address lies at the intersection: \emph{no existing method provides conformal coverage guarantees for extreme quantiles under mixing dependence}. Financial tail risk requires both EVT extrapolation and dependence-aware calibration.

% Paragraph 3: Key insight — combining GPD with switch coefficients
Our key insight is that GPD tail extrapolation and switch coefficient bounds compose naturally. The switch coefficient framework bounds coverage degradation as a function of the conformity scores' dependence structure, not the underlying data process. If we apply GPD extrapolation to conformity scores—transforming the calibration problem from estimating extreme empirical quantiles to estimating GPD parameters—the coverage guarantee transfers provided the extrapolated quantile estimate is conservative. The GPD fitting uses only exceedances above a moderate threshold, where empirical estimation is feasible; the switch coefficient bound applies to the score sequence, which inherits $\beta$-mixing from the return process. The combined guarantee provides finite-width prediction intervals at 99.9\%+ confidence with coverage gap bounded by $O(\beta(\tau) + n_u^{-1/2})$, where $\beta(\tau)$ is the mixing coefficient and $n_u$ is the number of threshold exceedances. For high-frequency financial data with thousands of daily observations, both terms are small.

% Paragraph 4: Additional challenge — circuit breaker censoring
A practical complication distinguishes financial tail events from environmental extremes. Exchange circuit breakers halt trading when prices move beyond specified limits—censoring the true extremes we seek to model. The May 6, 2010 flash crash saw the Dow Jones Industrial Average drop 1,000 points in minutes before trading halts intervened; the true counterfactual path remains unobserved. This censoring biases naive tail estimates downward: the largest observed returns understate the largest potential returns. We address this via a truncation correction to the GPD likelihood, treating circuit breaker activations as censored observations rather than true extremes. This correction is essential for any EVT analysis of financial data where trading halts occur, yet prior work on conformal prediction has not addressed it.

% Paragraph 5: Contributions
We make four contributions. First, we develop \emph{extreme conformal prediction} for financial time series, combining GPD tail extrapolation with the switch coefficient framework to provide coverage guarantees at extreme quantiles under $\beta$-mixing. Our coverage bound (Theorem~\ref{thm:coverage}) decomposes into mixing error $O(\beta(\tau))$ and estimation error $O(n_u^{-1/2})$, both of which are explicit and computable from data. Second, we introduce a censoring correction for circuit breaker effects, extending the GPD likelihood to account for truncated extremes. Third, we provide an impossibility result (Proposition~\ref{prop:impossibility}): when tail behavior does not vary with volatility regime—when returns are homoscedastic in the tails—GPD extrapolation cannot improve on empirical quantiles. This delimits when our method applies. Fourth, we validate the approach on two empirically relevant settings: FX markets with documented flash crashes (SNB peg removal 2015, Brexit 2016, JPY flash crash 2019) and cryptocurrency markets with extreme volatility (COVID crash March 2020, May 2021 deleveraging, FTX collapse November 2022). Across both settings, extreme conformal prediction achieves the nominal coverage while providing intervals 30--50\% narrower than those from classical conformal methods that must span the entire observed range.

\paragraph{Related work.}
The methodological foundation for this work is extreme value theory, specifically the peaks-over-threshold approach formalized by \citet{pickands1975statistical} and \citet{balkema1974residual}. For financial applications, \citet{mcneil2000estimation} establish GPD-based VaR estimation, and \citet{embrechts2013modelling} provide comprehensive treatment of EVT in quantitative risk management. Our use of GPD for conformity score extrapolation draws on \citet{pasche2025extreme}, who apply this idea to environmental extremes under approximate exchangeability; we extend their framework to dependent financial data.

Conformal prediction under dependence is the subject of active research. The switch coefficient framework of \citet{barber2025predictive} provides the theoretical foundation we build on, characterizing coverage degradation for $\beta$-mixing processes. Prior work by \citet{oliveira2024conformal} established $O(\sqrt{\tau_{\mathrm{mix}}/n})$ coverage bounds under mixing, which \citet{barber2025predictive} sharpen to $O((\ln n)/n)$ under geometric mixing. Online adaptive methods include Adaptive Conformal Inference \citep{gibbs2021adaptive,zaffran2022adaptive}, which adjusts thresholds based on recent miscoverage, and Temporal Conformal Prediction \citep{aich2025tcp}, which combines rolling calibration with quantile regression. These methods handle non-stationarity but do not address extreme quantile estimation.

For financial risk applications, \citet{schmitt2026taming} develop regime-weighted conformal calibration for portfolio VaR, using regime-similarity kernels to adapt to market conditions. Their approach works well for moderate quantiles (95--99\%) but relies on empirical quantiles, inheriting the infinite-interval problem at extreme levels. \citet{kato2024conformal} study conformal predictive portfolio selection, incorporating tail risk through interval optimization. Neither addresses EVT-based extrapolation.

The circuit breaker censoring problem has been studied in the market microstructure literature \citep{subrahmanyam1994circuit,kim2008price} but not in the context of conformal prediction or EVT estimation. Our truncation-corrected GPD likelihood adapts standard survival analysis techniques \citep{klein2003survival} to the financial setting.

% Paragraph 6: Paper structure
The remainder of the paper is organized as follows. Section~\ref{sec:background} reviews conformal prediction under dependence and the GPD tail approximation. Section~\ref{sec:method} develops extreme conformal prediction, including the censoring correction. Section~\ref{sec:theory} presents the coverage guarantee and impossibility result. Section~\ref{sec:simulations} evaluates the method through simulation studies calibrated to realistic financial dependence. Section~\ref{sec:empirical} applies the method to FX and cryptocurrency data spanning documented tail events. Section~\ref{sec:discussion} concludes with limitations and directions for future work.

% \input{sections/02_background}
% \input{sections/03_method}
% \input{sections/04_theory}
% \input{sections/05_simulations}
% \input{sections/06_empirical}
% \input{sections/07_discussion}

\bibliographystyle{plainnat}
\bibliography{references}

\end{document}
